%%
%% fall-lecture06.tex
%% 
%% Made by Alex Nelson <pqnelson@gmail.com>
%% Login   <alex@lisp>
%% 
%% Started on  2026-10-15T11:55:28-0700
%% Last update 2026-10-15T11:55:28-0700
%% 

\lecture{}

% This was the end of Lecture 5, but the proof began in Lecture 6
\begin{lemma}
Let $\xi$ be an $\RR^{n}$ bundle over a \emph{compact} base $B$. Then
for $k$ large enough, we can find a map $f\colon B\to G_{n}(\RR^{n+k})$
such that $\xi\iso f^{*}\gamma^{n}(\RR^{n+k})$.
\end{lemma}

\begin{proof}
The first key observation: it suffices to construct a continuous map
$\widehat{f}\colon E(\xi)\to\RR^{n+k}$ such that $\widehat{f}$ is
linear and injective on each fiber (since $\RR^{n}$ is the fiber of
$\xi$, $\RR^{n}$ is mapped to $\RR^{n}\subset\RR^{n+k}$). Then
construct $f\colon B\to G_{n}(\RR^{n+k})$ such that for all $b\in B$
we have $f(b)=\Im(\widehat{f}|_{\pi^{-1}(b)})$ is the image of
$\widehat{f}$ restricted to the fiber $\pi^{-1}(\{b\})$, and we can
check $f$ is continuous. We see
\begin{equation}
\vcenter{\xymatrix{%
\xi\ar[d] & \gamma^{n}\ar[d] & & e\ar@{|->}[d] & (f(\pi(e)),\widehat{f}(e))\ar@{|->}[d]\\
B\ar[r]^{f} & G_{n}(\RR^{n+k}) & & \pi(e)\ar@{|->}[r] & f(\pi(e))}}
\end{equation}
which suggests $e\mapsto(f(\pi(e)),\widehat{f}(e))$ is the map we are
trying to construct. The requirement is that $\widehat{f}(e)\in
f(\pi(e))$ to get a bundle map, since if $\pi(e)=b$ then $f\pi(e))=\Im(\widehat{f}|_{\pi^{-1}(b)})$.
Then we can use our previous lemma and find $\xi\iso f^{*}\gamma^{n}$.
We just deferred the hardest work (constructing $\widehat{f}$).

\textsc{Construct $\widehat{f}$.} We assumed $B$ is compact, so we can
find an open cover $\{U_{1},\dots,U_{r}\}$ of $B$ such that
$\xi|_{U_{i}}$ is trivial for each $i$. We can verify $B$ is normal,
so we can find $\{V_{1},\dots,V_{r}\}$ which is an open cover and for
all $i$, $\closure{V}_{i}\subset U_{i}$., Then we can find another
open cover $\{W_{1},\dots,W_{r}\}$ such that for all $i$ we have
$\closure{W}_{i}\subset V_{i}$. Then using Urysohn's lemma,
we can choose cutoff functions
\begin{equation}
\lambda_{i}\colon B\to\RR
\end{equation}
such that $\lambda_{i}|_{\closure{W}_{i}}=1$ (it is 1 on $\closure{W}_{i}$) and $\lambda_{i}|_{V_{i}^{\complement}}=0$
(it vanishes outside of $V_{i}$). We have the local trivialization
\begin{equation}
h_{i}\colon U_{i}\times\RR^{n}\to\pi^{-1}(U_{i}).
\end{equation}
We can look at some point in $\pi^{-1}(U_{i})$, then map into
\begin{align}
g_{i}\colon & \pi^{-1}(U_{i})\xrightarrow{h_{i}^{-1}}U_{i}\times\RR^{n}\xrightarrow{p}\RR^{n}\\
& h_{i}(b,v)\mapsto (b,v)\mapsto v.\nonumber
\end{align}
We can look at the fiber and find $g_{i}$ is linear and injective on
each fiber; i.e., $g_{i}(b,-)$ is injective and linear for each $b\in B$.
Now we want to go from local to global using the cut-off functions $\lambda_{i}$.
We define
\begin{equation}
g'_{i}(e)=\begin{cases}0 & \mbox{if }\pi(e)\notin U_{i}\\
\lambda_{i}(\pi(e))g_{i}(e) & \mbox{if }\pi(e)\in U_{i}.
\end{cases}
\end{equation}
The key thing is that $g'_{i}$ is continuous. Use
$\Support(\lambda_{i})\propersubset\closure{V}_{i}\subset U_{i}$ and
$g'_{i}$ is bounded, so it does not blow up on $\boundary V_{i}$.
Then $\widehat{f}\colon E(\xi)\to\RR^{nr}$ defined by
$f(e)=(g'_{1}(e), g'_{2}(e),\dots, g'_{r}(e))$ since we do not have
control on the overlaps we put the $g'_{i}$ into different components
of an $r$-tuple. Look at the base point $\pi(e)=b$, which belongs to
some $b\in U_{i}$, and this means $g'_{i}$ is linear and
injective. This concludes the proof.
\end{proof}

\begin{definition}
We define, as a set,
\begin{equation*}
\RR^{\infty}=\{(x_{1},x_{2},\dots,x_{n},\dots)\mid\mbox{only at most finitely many entries are nonzero}\}
\end{equation*}
and we can give it the \define{Direct Limit Topology} by considering
$\RR^{1}\subset\RR^{2}\subset\dots$, then
\begin{equation*}
\varinjlim_{i}\RR^{i}=\RR^{\infty},
\end{equation*}
and $U\subset\RR^{\infty}$ is open if and only if for all $i\in\NN$ we
have $U\cap\RR^{i}$ is open in $\RR^{i}$.
\end{definition}

\begin{definition}
We define the \define{Infinite Grassmann Manifold} $G_{n} := G_{n}(\RR^{\infty})$
as the set $G_{n}=\{\mbox{all $n$-dimensional subspaces of }\RR^{\infty}\}$
equipped with the topology obtained from the sequence
$G_{n}(\RR^{n})\subset G_{n}(\RR^{n+1})\subset\dots\subset G_{n}(\RR^{n+k})\subset\dots$
taking the directed limit $\varinjlim_{k}G_{n}(\RR^{n+k})$ and using
the direct limit topology: $U\subset G_{n}$ is open iff for all $k$ we
have $U\cap G_{n}(\RR^{n+k})$ is open.
\end{definition}

\begin{remark}
As a special case when $n=1$, we see
$G_{1}(\RR^{\infty})=\RP^{\infty}$ is the infinite-dimensional real
projective space, obtained by taking the direct limit of $\RP^{1}\subset\RP^{2}\subset\dots$.
\end{remark}

\begin{definition}
The \define{Universal Bundle} $\gamma^{n}$ over $G_{n}$ is the
$\RR^{n}$-bundle over $G_{n}$ consisting of
\begin{equation*}
\gamma^{n}=\{(X,v)\in G_{n}\times\RR^{\infty}\mid v\in X\}.
\end{equation*}
This is more commonly called the \define{Tautological Bundle}.
\end{definition}

\begin{lemma}
$\gamma^{n}$ is a bundle (i.e., satisfies the local triviality condition).
\end{lemma}

\begin{proof}
The proof is similar to the $\gamma^{n}(\RR^{n+k})$ proof, except
instead of $h\colon U\times\RR^{n}\to\pi^{-1}(U)$, we have a bunch of
these maps $h_{i}\colon U_{i}\times\RR^{n+i}\to\pi^{-1}(U_{i})$ for $i\in\NN$.
Then we just take the direct limit
$\varinjlim_{k}(U_{k}\times\RR^{n+i})$ and the key step is to check they
are compatible.
\end{proof}

\begin{remark}
Milnor and Stasheff~\cite[Lem~5.4]{milnor1974characteristic} offers a
different proof of this lemma.
\end{remark}

\begin{lemma}[Universal property]
Given an $\RR^{n}$-bundle $\xi$ over a \emph{paracompact} base $B$,
there exists a map $f\colon B\to G$ such that $f^{*}\gamma\iso\xi$.
\end{lemma}

\subsection{Notes to self}

I need to review paracompact spaces and their properties, since it has
been a while since I looked at them.

\begin{definition}
Recall, a topological space $X$ is \define{Paracompact} if $X$ is a
Hausdorff space and for every open cover $\{U_{\alpha}\}$ of $X$ there
exists an open cover $\{V_{\beta}\}$ such that
\begin{enumerate}
\item $\{V_{\beta}\}$ is a \emph{refinement} of $\{U_{\alpha}\}$: each
  $V_{\beta}$ is contained in some $U_{\alpha}$; and
\item $\{V_{\beta}\}$ is \emph{locally finite}: for any point $x\in X$
  there is a neighborhood $N\ni x$ which intersects with only finitely
  many of the $V_{\beta}$.
\end{enumerate}
\end{definition}

\begin{node}
Almost all familiar topological spaces are
paracompact. Second-countable spaces are paracompact.
\end{node}

\begin{theorem}[Stone]
Every metric space is paracompact.
\end{theorem}

\begin{theorem}[Morita]
If a regular topological space is the countable union of compact subsets,
then it is paracompact.
\end{theorem}

\begin{corollary}
The direct limit $K_{1}\subset K_{2}\subset\dots$ of compact spaces is paracompact.
\end{corollary}

\begin{theorem}[Dieudonn\'{e}]
Every paracompact space is normal.
\end{theorem}